% !TeX root = report.tex
\newcommand*{\PathToAssets}{../assets}%
\newcommand*{\PathToOutput}{../_output}%

%%%%%%%%%%%%%%%%%%%%%%%%%%%%%%%%%%%%%%
%% Compiled with XeLaTeX.
%%%%%%%%%%%%%%%%%%%%%%%%%%%%%%%%%%%%%%
\documentclass[12pt]{article}
\usepackage{my_article_header}
\usepackage{my_common_header}
\usepackage{float}  % [H] float placement: keep each exhibit in its own subsection

% Visible placeholder for an exhibit that the pipeline will emit but that is not
% built yet. Keeps the document compiling and the section outline complete.
\newcommand{\pending}[1]{%
  \begin{center}\fbox{\parbox{0.85\linewidth}{\centering\itshape
  [Exhibit pending: #1]}}\end{center}%
}

\begin{document}
\title{
Forecasting Crashes with a Smile\\
{\large A Replication and Extension of Martin \& Shi (2025)}
}
\author{Riley Haas \and Marija Jovicic}
\date{\today}

\begin{titlepage}
\maketitle

\doublespacing
\begin{abstract}
We replicate Martin \& Shi (2025), which shows that the shape of the
option-implied volatility smile carries a model-free signal for the physical
probability that a stock crashes. From index and single-name options we recover
risk-neutral return distributions, apply a power-utility fear correction, and
form Fr\'echet--Hoeffding bounds on the crash probability; the lower bound is
predicted to approximately equal the truth. Over the paper's 1996--2022 sample we
reproduce Table~1 to within a tight tolerance: the lower bound tracks the realized
crash frequency while the risk-neutral probability and the upper bound overstate
crash risk. We then extend every exhibit through 2025 and add sector-level crash
series and a Silicon Valley Bank case study. The full analysis regenerates
end-to-end from WRDS via a single \texttt{doit} command.
\end{abstract}
\end{titlepage}

\singlespacing

%======================================================================
\section{Introduction}
%======================================================================
Martin \& Shi (2025) derive option-implied \emph{bounds} on the physical
probability that an individual stock crashes over horizons of one to twelve
months. The risk-neutral crash probability, read directly from option prices,
``cries wolf'' --- it overstates crash risk, most severely in crises, because it
bundles a fear premium. The paper's contribution is a theory-based correction:
weighting by a power-utility investor's marginal value of a dollar and closing the
unknown stock--market dependence with copula bounds yields a \emph{lower bound}
that, under mild comonotonicity, should sit close to the true crash probability.

This project reproduces that result as a fully reproducible analytical pipeline.
Starting from raw CRSP and OptionMetrics data on WRDS, we build the cleaned
firm-month panel, recover the risk-neutral densities, form the bounds, and test
them against realized crashes. We then push every exhibit past the paper's 2022
cutoff and add extensions of our own.

\paragraph{How it went --- successes.} The core replication succeeds. Our
Table~1 matches the paper's cross-sectional statistics to within $0.007$: the
lower bound hugs the realized crash frequency while the risk-neutral probability
and (more so) the upper bound overstate crashes, exactly the pattern the theory
predicts. The entire pipeline --- pulls, cleaning, estimation, and exhibits ---
runs end-to-end from a clean checkout via \texttt{doit}, and every replicated
number is pinned to the paper by a tolerance unit test.

\paragraph{How it went --- challenges.} The harder problems were practical and
numerical rather than conceptual. (i) \emph{Numerical consistency}: the fitted
risk-neutral marginals (isotonic-regressed CDFs) and the option-price moment
integrals are two representations of the same distribution that can disagree at
short maturities, which briefly broke the theoretical bound ordering; we resolved
it by taking the boundary derivative in Result~5 from the fitted marginal, exactly
as Appendix~D prescribes. (ii) \emph{Scale}: forming a risk-neutral density for
every name, month, and maturity across a 35-million-row surface exhausted memory
until the pipeline was restructured to stream one day at a time. (iii) \emph{Data
gaps}: OptionMetrics is missing the July~2020 surface snapshot entirely, which we
handle by falling back to that month's last valid trading day. (iv)
\emph{Universe differences}: our CRSP membership vintage yields a modestly larger
firm set than the paper's, a difference we document rather than force.

%======================================================================
\section{Data}
%======================================================================
All raw data comes from WRDS and regenerates from scratch via \texttt{doit}; the
proprietary pulls cache locally and are never committed. The sample follows the
paper --- S\&P~500 index members observed on the last trading day of each month,
January 1996 through December 2022 --- and our extension carries the panel to the
most recent available data.

\begin{itemize}
  \item \textbf{OptionMetrics IvyDB} --- the standardized volatility surface
    (\texttt{vsurfd}) is the core input (the ``smile''); the zero-coupon curve
    (\texttt{zerocd}) supplies risk-free discounting; the security-name history
    (\texttt{secnmd}) resolves \texttt{secid} labels.
  \item \textbf{CRSP} --- the monthly stock file (\texttt{msf\_v2}, CIZ~2.0)
    provides realized returns and crash events with delisting returns folded into
    \texttt{mthret} (so wipeouts such as SVB register as crashes); the S\&P~500
    constituent list (\texttt{msp500list}) drives sample selection; the index file
    (\texttt{msix}) supplies the S\&P~500 level as the index spot.
  \item \textbf{WRDS CRSP--OptionMetrics link} joins option data (\texttt{secid})
    to CRSP returns (\texttt{permno}).
\end{itemize}

Financial firms are deliberately retained --- the universe is pure S\&P~500
membership with no industry screen, which is central to the SVB extension.

%======================================================================
\section{Methodology}
%======================================================================
For each firm and month we turn the option-implied smile into a crash-probability
estimate through five steps:

\begin{enumerate}
  \item \emph{Risk-neutral density.} Price out-of-the-money options on a fine
    moneyness grid, differentiate (Breeden--Litzenberger) to recover the
    risk-neutral CDF of the gross return, and enforce a valid distribution with an
    isotonic fit.
  \item \emph{Risk-neutral crash probability.} Read $P^{*}[R \le q]$ directly from
    that CDF.
  \item \emph{Fear correction.} Reweight by a power-utility investor's marginal
    utility ($\gamma = 2$, calibrated by the paper on the market) to move toward
    the physical measure.
  \item \emph{Copula bounds.} Close the unknown stock--market dependence with
    Fr\'echet--Hoeffding bounds; the comonotone \emph{lower bound} is the
    forecaster.
  \item \emph{Calibration test.} Regress the realized crash indicator on each
    probability --- the theory predicts a slope of one for the lower bound, not zero.
\end{enumerate}

We follow the paper's Appendix~D throughout: a fine grid $K/S \in [1/L, L]$
($L = 3$ at one to six months, $5$ at twelve), implied volatility interpolated
linearly within observed strikes and held flat outside, zero-dividend
Black--Scholes prices, and moments computed from option prices by the midpoint
rule. Table schemas are fixed centrally and every intermediate table is a tidy,
documented dataframe.

%======================================================================
\section{Replication (1996--2022)}
%======================================================================
We reproduce the paper's core exhibits over its original sample, in the order a
reader meets the argument: the raw signal for individual names and for the market
(Figures~\ref{fig:fig1} and~\ref{fig:fig2}), then the pooled summary statistics
(Table~\ref{tab:table1}) and the formal calibration test (Table~\ref{tab:table2}).

\subsection{Figure 1: single-name bounds}
The signal starts at the single-name level. Figure~\ref{fig:fig1} shows the
one-month bounds on a $20\%$ crash for two representative names; for each, the shaded
band between the ``cries wolf'' risk-neutral reading and the fear-corrected lower
bound is the copula uncertainty about the name's dependence on the market.

\begin{figure}[H]
\centering
\includegraphics[width=\linewidth]{\PathToOutput/fig1_single_name_bounds.png}
\caption{Bounds vary sharply across firms and over time: AIG's crash risk spikes in
the 2008--09 financial crisis, while Apple's is elevated in the early-2000s tech
unwind.}
\label{fig:fig1}
\end{figure}

\subsection{Figure 2: the market versus the cross-section}
Aggregating across names, Figure~\ref{fig:fig2} compares the cross-sectional median
of the bounds to the S\&P~500 index's own crash probability --- the $i = m$ case of
Result~3, computed from index options alone.

\begin{figure}[H]
\centering
\includegraphics[width=\linewidth]{\PathToOutput/fig2_median_bounds_market.png}
\caption{The market crash probability is consistently lower and smoother than the
typical constituent's --- the diversified index is less crash-prone than the stocks
that compose it.}
\label{fig:fig2}
\end{figure}

\subsection{Table 1: summary statistics}
Table~\ref{tab:table1} pools the bounds into summary statistics for every crash size
$q$ and horizon $\tau$, averaged both across firms (a monthly cross-section) and
across time (per firm).

\begin{table}[H]
\centering
\small
\caption{The mean lower bound sits next to the mean realized crash frequency, while
the risk-neutral probability and the upper bound run well above both --- the paper's
central finding. Our cross-sectional means match the paper to within $0.007$; $T$ is
the number of formation months, $N$ the number of firms.}
\label{tab:table1}
\input{\PathToOutput/table1.tex}
\end{table}

\subsection{Table 2: calibration regressions}
Table~\ref{tab:table2} runs the formal calibration test: for each measure $X$, the
pooled regression $I(R_{i,t\to t+\tau} \le q) = \alpha + \beta X + \varepsilon$ over
S\&P~500 constituents. A well-calibrated forecast has $\alpha \approx 0$ and
$\beta \approx 1$.

\begin{table}[H]
\centering
\footnotesize
\caption{The lower bound calibrates at $\beta \approx 1$, while the risk-neutral
probability and the upper bound are badly miscalibrated (slopes well below one) ---
the formal counterpart of Table~\ref{tab:table1}. Two-way clustered standard errors
(Thompson 2011) in parentheses; block-bootstrap standard errors (Martin \& Wagner
2019, 2500 samples) in brackets.}
\label{tab:table2}
\resizebox{\textwidth}{!}{\input{\PathToOutput/table2.tex}}
\end{table}

%======================================================================
\section{Extension Through 2025}
%======================================================================
We recompute every replicated exhibit over the sample extended past the paper's
December~2022 cutoff --- the first out-of-sample test of whether the pattern
survives. The genuinely new period is 2023 onward; its one notable crash event is
the March-2023 regional-bank stress (SVB and peers), while markets were otherwise
calm. Index membership past the CRSP vintage (2024-12) is held forward (the
\texttt{carried\_forward} freeze), which affects only the recent tail.

\subsection{Figures 1 and 2, extended}
Figures~\ref{fig:fig1ext} and~\ref{fig:fig2ext} carry the single-name and
cross-sectional bounds through the extended sample.

\begin{figure}[H]
\centering
\includegraphics[width=\linewidth]{\PathToOutput/fig1_single_name_bounds_ext.png}
\caption{Single-name bounds over the extended sample. The 2008 and 2020 spikes fall
inside the paper's own window; the genuinely new post-2022 tail is quiet.}
\label{fig:fig1ext}
\end{figure}

\noindent\textbf{What the new dates show.} The post-2022 period is calm for both
names: the one-month, $20\%$-crash upper bound peaks at only $2.9\%$ for Apple
(January 2023) and $2.1\%$ for AIG (May 2023, around the bank stress) --- an order of
magnitude below their 2008 and 2020 crisis readings. No new crisis-scale spike
appears, and the bounds behave out of sample exactly as they did within it.

\begin{figure}[H]
\centering
\includegraphics[width=\linewidth]{\PathToOutput/fig2_median_bounds_market_ext.png}
\caption{Cross-sectional median bounds and the index over the extended sample. The
market line stays lower and smoother than the median constituent, as in the
replication.}
\label{fig:fig2ext}
\end{figure}

\noindent\textbf{What the new dates show.} The aggregate picture is likewise quiet
after 2022: the cross-sectional median one-month upper bound stays below $1\%$
through 2023--2024, with the March-2023 bank stress registering as a small, localized
bump rather than a broad spike. The market-below-constituents ordering persists.

\subsection{Table 1, extended}
Table~\ref{tab:table1ext} recomputes the summary statistics over the full extended
sample.

\begin{table}[H]
\centering
\small
\caption{Summary statistics over the extended sample. The lower bound continues to
sit next to the realized crash frequency, while the risk-neutral probability and the
upper bound still overstate --- the replication's pattern is unchanged.}
\label{tab:table1ext}
\input{\PathToOutput/table1_ext.tex}
\end{table}

\noindent\textbf{What changed.} Almost nothing. Adding the post-2022 data moves every
summary statistic by at most $0.005$: at the $20\%$/12-month cell (across firms) the
realized frequency edges from $0.157$ to $0.153$, the lower bound from $0.121$ to
$0.120$, the risk-neutral probability from $0.235$ to $0.233$, and the upper bound
from $0.339$ to $0.335$. The lower-bound-tracks-realized relationship holds, nudged
microscopically downward by the calm post-2022 period.

\subsection{Table 2, extended}
Table~\ref{tab:table2ext} re-runs the calibration regressions over the full extended
sample.

\begin{table}[H]
\centering
\footnotesize
\caption{Calibration regressions over the extended sample. The lower bound's
$\beta \approx 1$ calibration survives, while the risk-neutral probability and the
upper bound remain miscalibrated.}
\label{tab:table2ext}
\resizebox{\textwidth}{!}{\input{\PathToOutput/table2_ext.tex}}
\end{table}

\noindent\textbf{What changed.} No significant change. The slope estimates move by
about $+0.02$: at the $20\%$/12-month cell the lower bound goes from $\beta = 1.13$ to
$1.15$, the risk-neutral from $0.68$ to $0.70$, and the upper bound from $0.42$ to
$0.44$. The lower bound stays calibrated ($\beta \approx 1$), the risk-neutral and
upper bounds stay miscalibrated, and the ranking is untouched. This is expected:
roughly two additional --- and mostly quiet --- years add little information relative
to the 27-year base sample, so the estimates barely move.

\subsection{Figure 6 --- out-of-sample forecasting performance}
Figure~\ref{fig:fig6ext} is our replication of the paper's Figure~6, carried through
the extended sample. Each line is a \emph{running} out-of-sample $R^2$ --- a
cumulative score of a forecaster against a naive benchmark, each firm's own
historical crash rate --- so any line above zero is beating that benchmark. The four
columns are the same $20\%$ crash measured $1$, $3$, $6$, and $12$ months ahead.

The pink line is our lower bound, computed with no fitting ($\alpha = 0$,
$\beta = 1$). The blue lines are the risk-neutral probability, which is known to run
too high: the paper corrects it \emph{statistically}, regressing realized crashes on
it over a trailing window each period (the calibration regression of
Table~\ref{tab:table2}) and using the fitted slope and intercept to rescale it going
forward. The two rows are the two versions of
that correction --- an expanding window (Panel~A) and a three-year rolling window
(Panel~B).

\begin{figure}[H]
\centering
\includegraphics[width=\linewidth]{\PathToOutput/fig6_oos_r2_ext.png}
\caption{Out-of-sample $R^2$ over the extended sample. The lower bound (pink) is
highest at every horizon and stays above zero for more than twenty years, while the
adjusted risk-neutral forecast (dark) declines and, at the $6$- and $12$-month
horizons, falls below zero after 2008--2010.}
\label{fig:fig6ext}
\end{figure}

\noindent\textbf{What the figure shows.} The lower bound is on top at every horizon
and stays above zero throughout, despite carrying no free parameters. The adjusted
risk-neutral does worse and, in the bottom-right panels, falls \emph{below} zero after
2008 --- the statistical fix does not merely underperform the bound, it underperforms
a firm's own historical average. The mechanism is a lag: the three-year window stays
calibrated to the crisis for three years after it ends, so it keeps trusting the
pessimistic risk-neutral number and cries wolf through the entire recovery. The lower
bound has nothing to lag --- it carries no window and no fitted coefficients, and
simply re-reads current option prices each month and re-prices. That contrast, between
a statistical correction anchored to a trailing window and a theory-based measure that
responds to live prices, is the central lesson of the paper.

\noindent\textbf{What the extension tests.} At the end of the paper's window every
line has been drifting down for about a decade. Two readings fit: the post-2012 market
was simply calm, with few crashes on which a forecaster could prove itself; or the
edge is genuinely eroding. The paper's data ends before the two can be told apart.
Carrying the sample through 2025 separates them, because it reintroduces real crash
events (the COVID aftermath, the 2022 drawdown, and the March-2023 bank stress). The
score is informative here by construction: a cumulative out-of-sample $R^2$ rises only
when a forecaster does better on the newly added data than it did on the past, so the
direction each line moves as the sample extends is a direct readout of whether the
edge is holding.

\noindent\textbf{What the new dates show.} The lower bound's edge holds. Across the
post-2022 data its cumulative $R^2$ moves by at most half a percentage point at any
horizon --- the $12$-month score edges from $13.2\%$ to $12.7\%$, the $6$-month from
$8.9\%$ to $8.6\%$ --- and the ranking is untouched: the lower bound stays on top and
above zero at every horizon, while the rolling-window risk-neutral stays below zero at
$6$ and $12$ months. Because the statistic is cumulative over a $27$-year base, two
mostly-quiet years cannot move it far by construction; what matters is that the path
is flat rather than turning down again. The bound forecast crashes over 2023--2024
about as well as it had over its whole history --- the benign reading, not the
erosion one --- which is the due-diligence result anyone relying on the paper needs:
the measure survived the events that came after the study.

%======================================================================
\section{Summary Statistics of the Underlying Data}
%======================================================================
Beyond the replicated exhibits, we present our own summary of the firm-month panel
--- the coverage of the option surface, the cross-section of firms, and the
distribution of realized returns --- so the reader can judge the data the bounds
are built on. Both exhibits are computed directly from the two cleaned inputs
(\texttt{clean\_surface.parquet} and \texttt{realized\_returns.parquet}), restricted
to S\&P~500 constituent member-months, so they describe exactly the sample the
replication runs on. This is distinct from Table~\ref{tab:table1}, which reports the
paper's statistics on the \emph{derived} crash-probability bounds; here we profile
the \emph{raw inputs} those bounds are built from.

Table~\ref{tab:eda} tabulates coverage year by year. The panel is broad and
strikingly stable --- roughly $500$ constituent names carry an option surface in
every year from 1996 to 2022, at a steady $\sim\!120$--$133$ quotes per firm-month
--- so the bounds are estimated on consistent data rather than a coverage that drifts
with the sample. Median implied volatility tracks the volatility cycle (peaking in
2000, 2008, and 2020), and the last column shows the object being forecast: the
realized frequency of a 20\% crash over the next twelve months is low on average but
highly time-varying, spiking above $0.45$ for surfaces formed in 2007--2008 and
falling near zero in calm years. Crashes are rare events, which is why the
out-of-sample $R^2$s in Figure~\ref{fig:fig6ext} are single digits by construction.

\begin{table}[H]
\centering
\small
\caption{Per-year coverage of the underlying option/return panel (S\&P~500
constituent member-months, 1996--2022). Coverage is broad ($\sim\!500$ names) and
continuous, at a stable quote density; median implied volatility traces the
volatility cycle; and the final column --- the realized 12-month 20\% crash frequency
--- shows the forecast target is a rare, strongly time-varying event, concentrated in
crisis years. The takeaway: the bounds rest on a consistent, well-populated panel, and
the thing they predict is genuinely hard to predict.}
\label{tab:eda}
\input{\PathToOutput/eda_coverage.tex}
\end{table}

Figure~\ref{fig:eda} reads the same panel four ways. Panel~(a) confirms the
cross-section is continuous, with only a brief thinning around the March-2020
dislocation. Panel~(b) is the most important for trusting the bound: it shows where on
the surface we actually have quotes, and the out-of-the-money-put tail (moneyness
$K/S$ well below $1$) --- the region the lower bound integrates --- is well populated,
especially at the longer maturities the paper uses for multi-month horizons. Panel~(c)
is the ``smile'' of the paper's title: median implied volatility rises sharply for
downside strikes, steeply at short maturities and flatter at long ones, the option
market's price of crash protection that the bound converts into a probability.
Panel~(d) shows the realized-return distribution fanning out and growing a fat left
tail as the horizon lengthens from one to twelve months, so the mass crossing the 20\%
and 30\% crash lines --- the events in the last column of Table~\ref{tab:eda} --- is
what the bounds are asked to forecast.

\begin{figure}[H]
\centering
\includegraphics[width=\linewidth]{\PathToOutput/eda_panel.png}
\caption{The underlying option/return panel at a glance (S\&P~500 constituent
member-months, 1996--2022). \textbf{(a)} The number of constituents carrying an option
surface each month is stable near $500$, dipping only briefly in March 2020.
\textbf{(b)} Quote availability across the moneyness $\times$ maturity grid: the
downside-put tail that the lower bound integrates is well covered, more so at longer
maturities. \textbf{(c)} The implied-volatility smile by maturity --- the downside
skew the bound reads as crash risk. \textbf{(d)} The realized gross-return
distribution develops a fat left tail as the horizon grows, placing real mass beyond
the 20\% and 30\% crash thresholds. Together they show the inputs are well-populated
exactly where the method needs them, and that crashes are rare but non-trivial.}
\label{fig:eda}
\end{figure}

%======================================================================
\section{Extensions}
%======================================================================
\paragraph{Sector crash series.} Applying the bound to the eleven Select Sector
SPDR ETFs (plus KRE) yields a \emph{direct} per-sector crash-probability series,
in contrast to the paper's average-of-constituents proxy.
\pending{Per-sector crash-probability series}

\paragraph{Silicon Valley Bank case study.} We read one event at three levels ---
the failed name (SIVB), the regional-bank ETF (KRE), and the financial sector
(XLF) --- to show the bound responding to the March~2023 stress in real time.
\pending{SVB case study (SIVB $\to$ KRE $\to$ XLF)}

\paragraph{The \texttt{crashbounds} package.} We package the pipeline as a
pip-installable library so a reader can call for a name's bounds and fear gap
directly.

%======================================================================
\section{Discussion and Conclusion}
%======================================================================
The replication confirms the paper's central claim on our independently
constructed panel: the option-implied lower bound is a well-calibrated forecast of
crash risk, materially tighter than the risk-neutral probability it corrects. The
main frictions were engineering and data-quality issues --- numerical consistency
between the two representations of the risk-neutral distribution, the memory cost
of the full cross-section, an OptionMetrics data gap, and membership-vintage
differences --- each of which we resolved or documented explicitly. The low
regression $R^2$ (a few percent) is expected: crashes are rare binary events, and
the meaningful comparison is the lower bound against the risk-neutral probability
on the same sample, where the bound wins.

\end{document}
